\section{Terminal-Bench Science task design}
\label{sec:tbscience}

Terminal-Bench Science asks for real scientific computational workflows that
are scientifically grounded, objectively and deterministically verifiable,
and difficult for substantive reasons; it targets a release solve rate around
10--20\% \citep{tbscience2026announcement,tbscience2026rubric}.  The
contribution process requires a proposal followed by a containerized Harbor
task and pytest-based verification.  The announced pull-request deadline is
August 17, 2026 \citep{tbscience2026announcement}.

\subsection{Recommended official objective}

The official task should be:
\begin{quote}
\emph{Improve the provided frozen AESP-LOCAPPR solver on a hidden suite of
personalized PageRank instances.  The submitted solver must satisfy the same
certified accuracy condition on every instance and reduce aggregate
instrumented degree-weighted graph-access work by a prescribed margin.}
\end{quote}
The task should not officially require the agent to prove
\(\softO(1/(\sqrt\alpha\epsppr))\), nor should it require beating the
instance-wise minimum of every known solver.

This narrower target has four advantages.
\begin{enumerate}[leftmargin=2em]
\item \textbf{It is outcome-verifiable.}  Correctness and work are both
      deterministic numerical quantities.
\item \textbf{It is scientifically coherent.}  The task asks the agent to
      improve a published accelerated local solver at its documented
      limitation, rather than solve an unbounded literature-review problem.
\item \textbf{It is demonstrably solvable.}  The two-stage hybrid is a private
      reference approach that can establish a realistic improvement threshold.
\item \textbf{It leaves room for discovery.}  The agent may find switching,
      restart, adaptive inner accuracy, priority scheduling, guarded
      over-relaxation, or a different mechanism.
\end{enumerate}

\subsection{Why not ``beat the best of all solvers''?}

Let \(W_a(i)\) be the work of an agent on instance \(i\), and let
\begin{equation}
  W_{\rm oracle}(i)
  :=\min\{W_{\AESP}(i),W_{\omega_\star\text{-SOR}}(i),
            W_{\LOCCH}(i),W_{\LOCGD}(i),W_{\APPR}(i)\}.
    \label{eq:oracle-baseline}
\end{equation}
Requiring \(W_a(i)<W_{\rm oracle}(i)\) broadly or on every instance is a
poor primary specification.  It makes solvability depend on five tuned
implementations, encourages portfolio selection rather than algorithmic
improvement, and complicates the verifier without improving scientific
clarity.

The broader solver set remains essential on the \emph{author side}.  Before
submission, the task author should verify that no included standalone baseline
passes the chosen threshold by itself.  In particular, replacing AESP with
plain optimal SOR should not be a trivial passing solution.  The recommended
separation is therefore
\begin{equation}
  \boxed{
  \begin{aligned}
  &\text{agent-facing baseline: frozen AESP-LOCAPPR},\\
  &\text{author-side calibration: AESP, optimal SOR, LocCH, LocGD, APPR}.
  \end{aligned}}
  \label{eq:benchmark-separation}
\end{equation}

\subsection{Verifier and score}

Let \(\mathcal I_{\rm hidden}\) be a deterministic hidden suite.  Every valid
submission must satisfy, for all \(i\in\mathcal I_{\rm hidden}\),
\begin{equation}
   \norm{D_i^{-1}(\widehat\pi_i-\pi_i)}_\infty\leq\epsppr^{(i)}.
   \label{eq:benchmark-correctness}
\end{equation}
The graph must be accessed through an instrumented interface that charges every
adjacency-list scan.  Full-graph matrix construction is permitted only if the
same accesses are charged; otherwise the task would reward bypassing the
scientific work model.

Define the geometric-mean work ratio
\begin{equation}
  R_{\rm work}
  :=\exp\left(
       \frac{1}{N}\sum_{i=1}^N
       \log\frac{W_a(i)}{W_{\AESP}(i)}
     \right).
  \label{eq:geomean-work-ratio}
\end{equation}
A plausible initial pass gate is
\begin{equation}
   R_{\rm work}\leq0.80,
   \label{eq:initial-speed-gate}
\end{equation}
provided the private hybrid reference passes with substantial margin, for
example \(R_{\rm work}\leq0.70\).  The exact gate must be calibrated using
frontier agents and the final container.  Add a regression guard such as
\begin{equation}
  Q_{0.90}\!\left(
      \frac{W_a(i)}{W_{\AESP}(i)}
  \right)\leq1.25,
  \label{eq:regression-guard}
\end{equation}
so a solver cannot pass by being exceptionally fast on a few cases and
catastrophically slow on many others.

The verifier should also check:
\begin{itemize}[leftmargin=2em]
\item deterministic output and work counts across repeated runs;
\item no modification of the frozen baseline or instrumentation;
\item finite memory and time limits;
\item correct handling of signed residuals and all hidden parameter regimes;
\item metadata recording sufficient to reproduce any failure.
\end{itemize}

\subsection{Hidden-suite design}

The hidden suite should vary graph family, source degree, \(\alpha\), and
\(\epsppr\).  It must include cases where AESP's early acceleration is valuable
and cases where late AESP overhead creates room for a hybrid.  A fixed switch
iteration should not pass.  Recommended strata are:
\begin{enumerate}[leftmargin=2em]
\item real sparse graphs with low-, medium-, and high-degree sources;
\item several small \(\alpha\) values;
\item at least three final tolerances per graph family;
\item star, path, long-spider, and core--boundary synthetic cases;
\item held-out graphs not named in the task description;
\item a few cases where continuing AESP is preferable, to penalize blind early
      switching.
\end{enumerate}

\subsection{Solvability witness and disclosure}

The task author should maintain a private reference implementation of the
AESP-to-LocSOR hybrid and use it only to calibrate the pass threshold.  The
agent-visible prompt should not reveal the switching rule.  The proposal must
disclose that the task originates from the authors' AESP and locally evolving
set research, and that a candidate hybrid reference solution exists.  The
TB-Science announcement explicitly requires provenance and prior use to be
disclosed \citep{tbscience2026announcement}.

\subsection{Acceptance assessment}

Under the rubric, the task is a strong candidate if the following are shown
before proposal submission:
\begin{enumerate}[leftmargin=2em]
\item a deterministic verifier enforces both accuracy and exact graph-access
      work;
\item the private reference hybrid clears the gate with comfortable margin;
\item all standalone baselines fail the official gate on the hidden aggregate;
\item the suite is small enough for reliable container execution but diverse
      enough to defeat fixed tuning;
\item pilot frontier-agent runs give a low but nonzero solve rate;
\item the task description is concise and does not reveal the hybrid method.
\end{enumerate}
Without a reference speedup or with only one or two public instances, the task
would be vulnerable to rejection as insufficiently demonstrated or too easy
to overfit.  Requiring a universal new theorem would create the opposite
problem: the outcome would no longer be efficiently verifiable.  The frozen
AESP improvement target is the appropriate middle ground.
